\section{Axial-profile evaluation}
\label{sec:s:profiles}

\paragraph{Metric.} Axial Wasserstein-1 compares normalised non-negative column profiles \cite{peyre2019}. For column $c$, write $p_{c,\ell}=x_{c,\ell}/\sum_m x_{c,m}$ and $\hat p_{c,\ell}=\hat x_{c,\ell}/\sum_m\hat x_{c,m}$, with cumulative sums $F_{c,\ell}$ and $\hat F_{c,\ell}$. Its discrete value is $\Delta z\sum_\ell|F_{c,\ell}-\hat F_{c,\ell}|$. We average using reference column brightness over columns above 10\,\% of its maximum, with $\Delta z=0.64\um$. Numerical normalisation uses a $10^{-12}$ denominator stabiliser. For unit-mass profiles, W1 is at least the absolute centroid difference and is zero only when the profiles coincide. It evaluates depth distributions, not connectivity or total brightness.

\begin{figure}[t]
\centering
\includegraphics[width=\linewidth]{profile_metric.pdf}
\caption{\textbf{A centroid can be correct while the axial profile is wrong.} Synthetic illustration: two layers and one merged layer share a centroid. Their W1 is the area between the cumulative profiles. This panel explains the metric; it is not a reconstruction result.}
\label{fig:s:profilemetric}
\end{figure}

\paragraph{Re-evaluation.} W1 was added to the existing reconstructions without retraining the plate or denoisers. Re-evaluations use the recorded angular subsets, noise realisations and validation-selected parameters. The new W1 tables report centroid error, PSNR and foreground SSIM from those same reconstructed volumes. Earlier encoder-control tables retain their recorded centroid metrics and are labelled accordingly.

The initial layered-column study required agreement within 0.05\,dB PSNR and $0.005\um$ centroid error for every row. TV met this criterion; the PnP configurations did not, and their profile re-evaluations were retained as descriptive analyses. A later W1 sweep checked the first row of each configuration against tolerances of 0.3\,dB and $0.02\um$; all checked configurations passed. Subsequent rows were recorded without this gate. \Cref{tab:s:profilereproduction} gives the largest deviations from the earlier metrics. The new evaluations follow the same reconstruction protocols but are not bit-identical replays.

\begin{table}[t]
\centering\small
\coltab{\begin{tabular}{@{}lccc@{}}\toprule
group & rows & max. $|\Delta$PSNR$|$ & max. $|\Delta$centroid$|$\\\midrule
initial profiles & 2100 & 0.198 & 0.0204\\
main & 1020 & 0.091 & 0.0353\\
objectives & 864 & 0.259 & 0.0117\\
dose & 1800 & 0.306 & 0.0216\\
compression & 552 & 0.210 & 0.0061\\
tolerance & 1104 & 0.220 & 0.0252\\
diffusion & 240 & 0.000 & 0.0000\\
\bottomrule\end{tabular}
}
\caption{Re-evaluation differences from the earlier recorded reconstructions. PSNR differences are in dB and centroid differences in $\mu$m. The EX2 group includes the strictly checked TV and descriptive PnP evaluations. The reported maxima include rows after the first-row checks.}
\label{tab:s:profilereproduction}
\end{table}

\paragraph{Full test set.} \Cref{tab:s:profilequality,fig:s:profilecrops} report the 60-crop comparison. Draws are averaged per crop before averaging crops. Cell-bootstrap samples retain all crops from each selected cell; confidence intervals concern the paired relative change of the means. The second optical seed changes the sampling during plate optimisation, with the initial phase fixed. The W1 compression sweep uses this initial random plate; independent random-re-alisation W1 sweeps were not completed.

\begin{table*}[t]
\centering\small
\begin{tabular}{@{}llcccc@{}}\toprule
decoder & plate & W1 & centroid & PSNR & FG-SSIM\\\midrule
TV & random & 0.371 & 0.259 & 36.70 & 0.850\\
TV & learned s0 & 0.218 & 0.150 & 38.29 & 0.910\\
TV & learned s1 & 0.218 & 0.149 & 38.31 & 0.910\\
pretrained PnP & random & 0.579 & 0.490 & 35.72 & 0.806\\
pretrained PnP & learned s0 & 0.290 & 0.224 & 37.56 & 0.900\\
pretrained PnP & learned s1 & 0.289 & 0.222 & 37.58 & 0.900\\
neuron-trained PnP & random & 0.271 & 0.197 & 38.01 & 0.902\\
neuron-trained PnP & learned s0 & 0.195 & 0.139 & 39.60 & 0.934\\
neuron-trained PnP & learned s1 & 0.194 & 0.136 & 39.64 & 0.934\\
\bottomrule\end{tabular}

\caption{Reconstruction metrics from the same re-evaluations ($C=8$, 60 crops, five draws). W1 and centroid error are in $\mu$m; PSNR is in dB. Learned PnP uses the combined neuron-trained 2D and 3D denoising configuration.}
\label{tab:s:profilequality}
\end{table*}

\begin{figure*}[t]
\centering
\includegraphics[width=0.95\linewidth]{profile_per_crop.pdf}
\caption{\textbf{Per-crop axial-profile error.} Random versus learned plate with three decoders at $C=8$, after averaging five direction subsets. Filled markers: 12 prespecified crops; open markers: 48 additional random crops. Every crop has lower W1 with the learned plate.}
\label{fig:s:profilecrops}
\end{figure*}

\begin{figure}[t]
\centering
\includegraphics[width=\linewidth]{profile_recovery.pdf}
\caption{Summary of axial W1 under three reconstruction methods. The learned and random plates use all 60 crops; the DMD and depth-blind references use the 12 prespecified crops. Learned PnP denotes the neuron-trained configuration.}
\label{fig:s:profilesummary}
\end{figure}

\paragraph{Representative volume.} Before viewing the candidate images, we selected the crop closest to the median of the mean fractional W1 reduction across the three reconstruction methods. Each crop's error is first averaged over its five draws; ties are resolved by crop identifier. This selects SOM006\_cell2\_x7. The main paper shows pretrained PnP at draw 0; \cref{fig:s:neuronpriors} shows all three methods on the same draw. Lateral colour indicates the depth of the brightest voxel, modulated by maximum intensity with a shared power transform of 0.7. Axial views use a common linear grey scale. This example illustrates depth placement and profile recovery; it is not a neuronal-connectivity demonstration.

\begin{figure*}[t]
\centering
\includegraphics[width=\linewidth]{neuron_priors.pdf}
\caption{Axial maximum-intensity projections of the fixed representative crop across all three decoders. All tiles use the same reference-based intensity scale. Numbers are the full-volume W1 for this draw, not projection errors.}
\label{fig:s:neuronpriors}
\end{figure*}

\paragraph{Layered columns.} A signal column is layered when its reference profile has at least two peaks, each with height and prominence at least 20\,\% of the column maximum. Of the 60 crops, 59 contain such columns. Under the strict TV reproduction check, their W1 falls from $0.467$ to $0.309\um$, a 33.8\,\% reduction (95\,\% CI: 31.7--36.3\,\%). The relative reduction is not larger than the corresponding centroid improvement. Peak recovery remains incomplete: overall, the learned and random plates recover both outer layers in 12 versus 4\,\% of layered columns with TV, 21 versus 12\,\% with pretrained PnP, and 20 versus 13\,\% with learned PnP. These peak-based summaries are descriptive.

\section{Measurement distinguishability and crossing structures}
\label{sec:s:separation}

\paragraph{Merged-depth alternatives.} In each layered reference column, we move the fine-grid intensity to its axial centroid, distributing mass between the two neighbouring planes when needed. This preserves the lateral projection and centroid while changing the depth arrangement. Other columns remain unchanged. The comparison uses the fine-grid forward model, not its piecewise-constant reconstruction approximation. One crop has no layered columns and is excluded, leaving 59.

\paragraph{Poisson comparison.} Let $s$ and $s'$ be the expected signal-count vectors of the reference and alternative, and let $b$ be the known background. With $\mu=s+b$ and $\nu(\beta)=\beta s'+b$, we compute
\begin{equation}
\min_{\beta>0}\sum_i\left[
\mu_i\log\frac{\mu_i}{\nu_i(\beta)}+\nu_i(\beta)-\mu_i
\right].
\label{eq:s:poisson}
\end{equation}
At zero background the optimum is $\beta=\sum_i s_i/\sum_i s'_i$. If $p=s/\sum_i s_i$ and $q=s'/\sum_i s'_i$, the profiled divergence equals $(\sum_i s_i)D_{\mathrm{KL}}(p\Vert q)$. It therefore compares the expected count patterns after removing a global brightness difference. We average over the five angular subsets for each crop and report the geometric mean learned/random ratio with a cell-bootstrap interval. With background, brightness is fitted numerically to the signal while the background stays fixed.

The primary ratio is 10.0 [8.7, 11.8], and every eligible crop improves. Added background gives 10.4 [9.0, 12.4]. Random-plate divergences are already substantial for some alternatives; this result measures increased separation rather than proving that an alternative is measurement-consistent. Swapped-depth, mirrored-depth and equal-W1 rigid-shift alternatives on phantoms provide additional comparisons. The improvement is not uniformly larger for merging than for a rigid shift of matched transport distance, so we make no preferential structural-versus-centroid claim.

\paragraph{Crossing phantoms.} Two Gaussian filaments of width $\sigma=0.32\um$ cross at the lateral field centre. The three orientation/crossing-angle pairs are $(0^\circ,90^\circ)$, $(30^\circ,60^\circ)$ and $(60^\circ,45^\circ)$. Axial centres are reconstruction-plane indices 7 and 5; only the former coincides with the mirror pairing used by the plate objective. Gaps are 0, 2, 3, 4, 6 and 8 planes, with $0.64\um$ per plane. The zero-gap case is a merged control. We use five measurement draws, the same photon protocol as the neuronal volumes and the existing validation-selected reconstruction parameters.

\paragraph{Success rule.} Crossing columns contain both reference filaments and resolve their peaks. Estimated peaks must match the two reference depths within one plane, with height and prominence at least 20\,\% of the column maximum. At least half of the reference crossing-column brightness must satisfy this condition. Correct ordering additionally requires the estimated intensity centroids along the two separate filament arms to follow the reference depth order, with a margin greater than half a plane. Arm columns contain at least 20\,\% of their own filament's maximum projected intensity and less than 2\,\% of the other's. A volume succeeds only when both tests pass.

The reported threshold requires at least 80\,\% success among the 15 geometry/draw combinations at each axial centre, at that gap and every larger tested gap. TV and pretrained PnP improve from 6 to 4 planes; learned PnP remains at 6. Neither plate produces false splitting in the merged control. The main figure pools the two centres for its curves; the threshold is checked at each centre separately. Unequal-brightness filaments use a weaker branch of amplitude 0.3. Closely spaced three-filament crossings remain difficult, and the selected-alternative KL does not by itself guarantee their recovery. Pooled KL/success associations are confounded by gap and phantom type and are not used as a predictive guarantee.

\section{Profile error under acquisition and model changes}
\label{sec:s:profileconditions}

\begin{table}[t]
\centering\small
\coltab{\begin{tabular}{@{}llcccccc@{}}
\toprule
Reconstruction & Plate & W1, $C=1$ & Centroid & W1, $C=4$ & Centroid & W1, $C=8$ & Centroid\\
\midrule
TV & Random & 0.249 & 0.187 & 0.298 & 0.223 & 0.318 & 0.243\\
TV & Learned, seed 0 & 0.170 & 0.135 & 0.208 & 0.166 & 0.196 & 0.153\\
TV & Learned, seed 1 & 0.160 & 0.124 & 0.202 & 0.158 & 0.195 & 0.151\\
Pretrained PnP & Random & 0.413 & 0.357 & 0.445 & 0.384 & 0.496 & 0.431\\
Pretrained PnP & Learned, seed 0 & 0.222 & 0.184 & 0.254 & 0.212 & 0.277 & 0.231\\
Pretrained PnP & Learned, seed 1 & 0.213 & 0.173 & 0.249 & 0.206 & 0.279 & 0.234\\
Neuron-trained PnP & Random & 0.203 & 0.161 & 0.234 & 0.184 & 0.282 & 0.221\\
Neuron-trained PnP & Learned, seed 0 & 0.142 & 0.110 & 0.167 & 0.130 & 0.206 & 0.163\\
\bottomrule
\end{tabular}
}
\caption{Axial W1 ($\mu$m) versus compression on the 12 prespecified crops at equal total deposited dose. Fewer directions receive longer exposure per image.}
\label{tab:s:profilecompression}
\end{table}

\begin{figure}[t]
\centering
\includegraphics[width=\linewidth]{measurements.pdf}
\caption{Learned and random plates under fewer illumination directions. W1 is evaluated on the 12 prespecified crops. The horizontal reference is each decoder's random-plate result at $C=1$.}
\label{fig:s:profilecompression}
\end{figure}

\paragraph{Dose and noise.} \Cref{tab:s:profiledose,tab:s:profilenoise} give W1 over the 12 prespecified crops. Dose-dependent reconstruction parameters are selected on the validation volumes using the same grid-extension rule. At 0.3 times dose the learned plate improves on the random plate at full dose under all three methods. Dose levels 0.1 and 3 provide the remainder of the curve. A plateau of these curves is a property of the tested inference protocols, not proof that additional photons could never help another reconstruction method.

\begin{table}[t]
\centering\small
\coltab{\begin{tabular}{@{}lccccc@{}}
\toprule
Reconstruction & Nominal dose & W1, random & W1, learned & Centroid, random & Centroid, learned\\
\midrule
TV & 0.1 & 0.466 & 0.292 & 0.365 & 0.236\\
TV & 0.3 & 0.442 & 0.234 & 0.341 & 0.185\\
TV & 1 & 0.318 & 0.196 & 0.243 & 0.153\\
TV & 3 & 0.315 & 0.192 & 0.242 & 0.151\\
Pretrained PnP & 0.1 & 0.834 & 0.407 & 0.733 & 0.352\\
Pretrained PnP & 0.3 & 0.631 & 0.302 & 0.550 & 0.251\\
Pretrained PnP & 1 & 0.496 & 0.277 & 0.431 & 0.230\\
Pretrained PnP & 3 & 0.494 & 0.276 & 0.430 & 0.230\\
Neuron-trained PnP & 0.1 & 0.345 & 0.239 & 0.269 & 0.186\\
Neuron-trained PnP & 0.3 & 0.304 & 0.220 & 0.236 & 0.174\\
Neuron-trained PnP & 1 & 0.282 & 0.207 & 0.221 & 0.164\\
Neuron-trained PnP & 3 & 0.275 & 0.203 & 0.218 & 0.162\\
\bottomrule
\end{tabular}
}
\caption{Axial W1 ($\mu$m) versus deposited excitation dose at $C=8$. Parameters are reselected on validation data for each condition.}
\label{tab:s:profiledose}
\end{table}

\begin{table}[t]
\centering\small
\coltab{\begin{tabular}{@{}llcccccc@{}}
\toprule
Condition & Reconstruction & W1, random & W1, learned & Centroid, random & Centroid, learned & W1 change & Centroid change\\
\midrule
Shot noise & TV & 0.318 & 0.196 & 0.243 & 0.153 & $-38$\,\% [$-43$, $-34$] & $-37$\,\% [$-41$, $-34$]\\
Shot noise & Pretrained PnP & 0.496 & 0.277 & 0.431 & 0.231 & $-44$\,\% [$-50$, $-39$] & $-46$\,\% [$-53$, $-41$]\\
Shot noise & Neuron-trained PnP & 0.282 & 0.206 & 0.221 & 0.163 & $-27$\,\% [$-34$, $-20$] & $-26$\,\% [$-35$, $-18$]\\
Background + read noise & TV & 0.324 & 0.227 & 0.249 & 0.180 & $-30$\,\% [$-38$, $-22$] & $-28$\,\% [$-36$, $-20$]\\
Background + read noise & Pretrained PnP & 0.627 & 0.297 & 0.548 & 0.247 & $-53$\,\% [$-57$, $-48$] & $-55$\,\% [$-60$, $-50$]\\
Background + read noise & Neuron-trained PnP & 0.288 & 0.210 & 0.228 & 0.166 & $-27$\,\% [$-35$, $-19$] & $-27$\,\% [$-36$, $-18$]\\
0.1 dose & TV & 0.466 & 0.292 & 0.365 & 0.236 & $-37$\,\% [$-45$, $-29$] & $-35$\,\% [$-44$, $-28$]\\
0.1 dose & Pretrained PnP & 0.834 & 0.407 & 0.733 & 0.352 & $-51$\,\% [$-57$, $-45$] & $-52$\,\% [$-58$, $-45$]\\
0.1 dose & Neuron-trained PnP & 0.345 & 0.239 & 0.269 & 0.186 & $-31$\,\% [$-37$, $-24$] & $-31$\,\% [$-38$, $-24$]\\
\bottomrule
\end{tabular}
}
\caption{Axial W1 ($\mu$m) with background/read noise or reduced dose. Reductions compare paired means; intervals use cell resampling.}
\label{tab:s:profilenoise}
\end{table}

\paragraph{Optical objectives.} The W1 comparison across both seeds and all three decoders is in \cref{tab:s:profileobjectives}. It uses the same optical initialisation and optimisation budget. The penalty families and their implementation are described in \cref{sec:s:objectives}.

\begin{table}[t]
\centering\small
\coltab{\begin{tabular}{@{}llcc@{}}\toprule
comparison & decoder & W1 reduction [95\,\% CI] & improved\\\midrule
coherence, s0 & TV & $17$\,\% [$15$, $20$] & 12/12\\
coherence, s0 & pretrained PnP & $20$\,\% [$16$, $25$] & 12/12\\
coherence, s0 & learned PnP & $13$\,\% [$12$, $15$] & 12/12\\
coherence, s1 & TV & $16$\,\% [$14$, $19$] & 12/12\\
coherence, s1 & pretrained PnP & $17$\,\% [$13$, $23$] & 11/12\\
coherence, s1 & learned PnP & $13$\,\% [$11$, $15$] & 12/12\\
Fisher, s0 & TV & $14$\,\% [$12$, $16$] & 12/12\\
Fisher, s0 & pretrained PnP & $19$\,\% [$17$, $22$] & 12/12\\
Fisher, s0 & learned PnP & $14$\,\% [$12$, $17$] & 11/12\\
Fisher, s1 & TV & $14$\,\% [$12$, $16$] & 12/12\\
Fisher, s1 & pretrained PnP & $16$\,\% [$12$, $20$] & 12/12\\
Fisher, s1 & learned PnP & $16$\,\% [$14$, $18$] & 12/12\\
\bottomrule\end{tabular}
}
\caption{W1 reduction of our plate against each alternative objective ($C=8$, 12 prespecified crops). Seeds differ in optimisation sampling. Counts refer to crops with lower W1.}
\label{tab:s:profileobjectives}
\end{table}

\paragraph{Alignment.} \Cref{tab:s:profilealignment} complements the earlier centroid-error table with W1. Rebuilding the operator for the displaced plate retains the gain. Using the nominal operator under axial or lateral displacement eliminates it; the one-feature lateral shift makes the learned plate worse on W1. Pointing errors retain a smaller advantage. This comparison supports the need for calibration, not unconditional alignment robustness.

\begin{table}[t]
\centering\small
\coltab{\begin{tabular}{@{}lcccccc@{}}
\toprule
Condition & W1, random & W1, learned & Centroid, random & Centroid, learned & W1 change & Centroid change\\
\midrule
none (design) & 0.282 & 0.206 & 0.221 & 0.163 & $-27$\,\% [$-34$, $-20$] & $-26$\,\% [$-35$, $-18$]\\
Known: D=6.1 & 0.285 & 0.205 & 0.225 & 0.162 & $-28$\,\% [$-35$, $-22$] & $-28$\,\% [$-37$, $-20$]\\
Known: D=7.1 & 0.284 & 0.209 & 0.226 & 0.166 & $-27$\,\% [$-34$, $-20$] & $-26$\,\% [$-35$, $-19$]\\
Nominal operator: scan tilt 0.005 & 0.318 & 0.254 & 0.253 & 0.196 & $-20$\,\% [$-25$, $-16$] & $-22$\,\% [$-30$, $-15$]\\
Nominal operator: scan tilt 0.01 & 0.379 & 0.328 & 0.308 & 0.251 & $-13$\,\% [$-16$, $-10$] & $-18$\,\% [$-22$, $-15$]\\
Nominal operator: D=6.1 & 0.563 & 0.586 & 0.528 & 0.571 & $+4$\,\% [$+0.3$, $+8$] & $+8$\,\% [$+3$, $+14$]\\
Nominal operator: D=7.1 & 0.447 & 0.473 & 0.383 & 0.438 & $+6$\,\% [$+5$, $+7$] & $+14$\,\% [$+12$, $+17$]\\
Nominal operator: lateral shift 1 feature & 0.451 & 0.500 & 0.372 & 0.400 & $+11$\,\% [$+2$, $+23$] & $+7$\,\% [$-0.4$, $+19$]\\
\bottomrule
\end{tabular}
}
\caption{Alignment sensitivity in W1 ($\mu$m, learned PnP, $C=8$). Changes compare learned against random under the same perturbation. Updated-model rows rebuild the operator and reselect parameters on validation.}
\label{tab:s:profilealignment}
\end{table}

\FloatBarrier
